%% kapa: exam=1 round=37 year=2026 period=1 form=A subject=economics subject_ko=경제학원론 seq=31 number=71
%% source: Q-Net 감정평가사 기출문제 (cst00309) artlSeq=5258924 file=raw/1008/5258924/2026년 제37회 감정평가사 1차 1교시.pdf
%% conversion: kordoc 4.24.51 -> tools/gen1.py (auto); review=proofread-sonnet; flags=none
\begin{question}{71}
아래의 단기 필립스곡선에 관한 설명으로 옳지 않은 것은? (단, $\pi$, $\pi^{e}$, $u$, $u_{n}$은 각각 인플레이션율, 기대 인플레이션율, 실업률, 자연실업률이다.)
\begin{framedtext}
\centering $\pi-\pi^{e}=-2(u-u_{n})$
\end{framedtext}
\begin{choices}
  \item 총공급곡선이 우상향할 때 나타날 수 있는 관계이다.
  \item 총수요곡선의 위치가 이동할 때 나타날 수 있는 관계이다.
  \item 실업률을 1 \%p 낮추려면 인플레이션율을 2 \%p 상승시켜야 한다.
  \item 인플레이션율이 기대 인플레이션율보다 낮은 경우 고용은 완전고용 수준보다 높다.
  \item 기대 인플레이션율을 낮추면 높은 인플레이션율의 희생없이 실업률을 낮출 수 있다.
\end{choices}
\end{question}
\answer{4}
